\section{Root neighborhoods and a split residue map}\label{sec:covers}

We begin the proof of Theorem~\ref{thm:supported-lifting}.  A generated
multiple of the actual adjoint restriction constructs a morphism from
the reduced boundary to projective space.  Near its compact fibers we
will replace the supported divisor by a reduced Cartier covering divisor
and study sections on its finite neighborhoods.  The use of neighborhood
covers and successive finite thickenings has its threefold antecedents in
\cite[Sections~1--2 and~4]{Miyaoka1988} and
\cite[Section~4]{Kawamata1992}.  Here no extension of the boundary map
off the reduced boundary is assumed.

Fix the data of Theorem~\ref{thm:supported-lifting}, and put
\(S=\floor B\), so \(\Supp P=\Supp S\).  Choose a sufficiently divisible even integer
\(q>0\) which clears the actual adjoint and equivalence indices, the
boundary and \(P\) coefficients, and a generated degree of \(A|_S\).
There are then an actual line and section
\begin{equation}\label{eq:lifting-data}
 L=\OO_V(qA),\qquad s\in H^0(V,L),\qquad
 G=\operatorname{div}_L(s)=qP=\sum_i d_iS_i,
\end{equation}
where every \(d_i\) is a positive integer, and a morphism
\begin{equation}\label{eq:boundary-map}
 f:S\longrightarrow T=\PP^b,\qquad L|_S\simeq f^*\OO_T(1).
\end{equation}
The morphism is constructed from a generating system on \(S\); it is
used only on that reduced subspace and is not required to be surjective.
Choose \(r>0\) divisible by \(q\) and
every \(d_i\), and put \(a=r/q\).  In particular \(a\) is a positive
integer and \(d_i\leq r\).

\subsection{Root pairs near a compact analytic subspace}

\begin{lemma}[A root with a prescribed boundary comparison]
\label{lem:root-neighborhood}
Let \(A\) be a compact analytic subspace of a Hausdorff complex analytic
space \(X\), and let \(L\) be a holomorphic line on a neighborhood of
\(A\).  Suppose a line \(Q\) on \(A\) and an isomorphism
\(\beta:Q^{\otimes r}\simeq L|_A\) are given, with \(r>0\).
There are a neighborhood \(U\) of \(A\), a line \(P\) on \(U\), and
isomorphisms
\[
 \alpha:P^{\otimes r}\simeq L|_U,\qquad \gamma:P|_A\simeq Q,
 \qquad \alpha|_A=\beta\circ\gamma^{\otimes r}.
\]
If two root pairs are already defined near \(A\), every prescribed
power-compatible isomorphism between their restrictions to \(A\) extends
to an isomorphism on a neighborhood, uniquely as a germ about \(A\).
The root pair itself is not asserted unique.
\end{lemma}
\begin{proof}
The analytic Kummer sequence
\[
 1\longrightarrow\mu_r\longrightarrow\OO_X^\times
 \xrightarrow{(\cdot)^r}\OO_X^\times\longrightarrow1
\]
is exact also for singular or nonreduced analytic spaces.  A unit germ
has a local holomorphic root; the kernel is the locally constant sheaf
of \(r\)th roots of unity, since \(u^r-1\) has distinct roots, including
in a local ring with nilpotents.  Line bundles are classified by
\(H^1(\OO^\times)\), and abelian torsors by \(H^1\) of the corresponding
sheaf \cite[Tags~09NU and~02FQ]{StacksProject}.

For an abelian sheaf \(F\) and a compact subset whose distinct points
have disjoint neighborhoods, SGA~4 continuity
\cite[Expos\'e~Vbis, Lemma~4.1.3]{SGA4Continuity} gives
\begin{equation}\label{eq:compact-continuity}
 \varinjlim_{U\supset A}H^i(U,F)=H^i(A,F|_A)
 \quad\text{for every }i.
\end{equation}
The separation hypothesis holds in the Hausdorff analytic space.  Although
the intrinsic structure sheaf of \(A\) need not be the inverse image of
\(\OO_X\), the inverse image of \(\mu_r\) is exactly \(\mu_r\) on \(A\).
Naturality of Kummer therefore makes the obstruction to a root of \(L\)
restrict to zero in \(H^2(A,\mu_r)\), because \(Q\) is such a root there.
Continuity in degree two kills that obstruction on a smaller neighborhood,
giving a root pair \((P,\alpha)\).

The difference between \((P|_A,\alpha|_A)\) and \((Q,\beta)\) is a
\(\mu_r\)-torsor.  By continuity in degree one it extends to a torsor on
a smaller neighborhood.  The associated line with its trivialized
\(r\)th power twists \(P\) to give the prescribed pair on \(A\).
Finally, the sheaf of compatible isomorphisms of two root pairs is itself
a locally trivial finite \(\mu_r\)-torsor.  A prescribed section on \(A\)
trivializes its restriction there.  Continuity in degree one therefore
trivializes this torsor on a smaller neighborhood; choose a section there.
The ratio of its boundary restriction to the prescribed section is a
section of \(\mu_r\) on \(A\).  Continuity in degree zero extends that
ratio uniquely as a germ, and correcting the chosen section gives the
prescribed extension.  The same degree-zero injectivity gives uniqueness
as a germ.  These sections are holomorphic, being locally solutions of
\(z^r=u\) for a holomorphic unit.
\end{proof}

Apply the lemma near a parameter \(t\in T\) to
\(A=(f^{-1}(t))_{\red}\), using a local frame \(\ell\) of \(\OO_T(1)\)
and the trivial root whose \(r\)th power is \(f^*\ell\).  We obtain a
root pair \(P_1^r\simeq L\) on a neighborhood of that compact fiber.
The prescribed compatible frame on the fiber extends as a germ inside
\(S\), by the last part of the lemma.  If it is defined on
\(N_S\subset S\) containing the fiber, properness of \(f\) gives a
smaller base neighborhood \(U\ni t\) with \(S_U=f^{-1}(U)\subset N_S\):
remove the closed set \(f(S\setminus N_S)\).  Shrink once more so this
whole \(S_U\) lies in the ambient root neighborhood \(W_0\), and replace
\(W_0\) by \(W=W_0\setminus(S\setminus S_U)\).  Thus
\begin{equation}\label{eq:local-root}
 W\cap S=S_U,\qquad P_1^r\simeq L|_W,\qquad
 p_1\text{ a frame of }P_1|_{S_U},\quad p_1^r=f^*\ell.
\end{equation}
Comparisons between two choices, with prescribed power-compatible
boundary frames, extend uniquely as ambient germs near a compact fiber.
Only this torsion comparison is extended; no arbitrary boundary
trivialization is extended off \(S\).

\subsection{The first normalized cover and the lifting target}

On a chart \eqref{eq:local-root}, take the full normalized cyclic cover
\(\pi:Z\to W\) defined by
\[
 y^r=\pi^*s\quad\text{in }\pi^*P_1^r,
\]
retaining all components and the \(\mu_r\)-action.  Finite analytic
normalization is available by
\cite[Part~B, Section~4, Corollaries~2--3]{Houzel1961}.  At a general point of
\(S_i\), a root of a local unit reduces the equation to \(y^r=x^{d_i}\).
Because \(d_i\mid r\), each normalized branch is
\[
 x=t^{r/d_i},\qquad y=\zeta t.
\]
The ramification index is \(r/d_i\) and \(y\) has order one.  Outside
\(S\) the equation roots a unit and is \'etale.  A finite map preserves
the dimension of a prime analytic subset, so no further divisor above a
codimension-two set is missed.  The zero divisor
\[
 E=(y=0)
\]
is Cartier: \(y\) is a nonzerodivisor on each normal component.  A Cartier
divisor on a normal space is \(S_1\), and the calculation makes it
generically reduced; hence it is reduced.  Its canonical line identity is
\begin{equation}\label{eq:cover-line}
 \OO_Z(E)\simeq\pi^*P_1,
\end{equation}
sending the canonical section of \(\OO_Z(E)\) to \(y\).

Put
\begin{equation}\label{eq:cover-layers}
 I=\OO_Z(-E),\quad g=f\circ\pi|_E:E\to U,\quad
 \mathcal A_j=I^j/I^{j+1}=\OO_E(-jE)\quad(j\in\Z).
\end{equation}
The map \(E\to S_U\) is finite and \(g\) is proper.  Each \(\mathcal A_j\)
is invertible on the possibly singular reduced \(E\).  In the Laurent
graded algebra of these layers, \(u=y/p_1\) denotes the degree-one frame.
This is intrinsic on the associated graded: locally divide the equation
of \(E\) by the boundary value of a frame of \(P_1\).  It does not assert
that \(p_1\) is an ambient frame.

For \(j\in\Z\) and \(k\geq1\), put
\begin{equation}\label{eq:truncations}
 \mathcal T_{j,k}=I^j/I^{j+k}.
\end{equation}
We view this as a sheaf of complex vector spaces on the underlying
topological space of \(E\).  A sheaf on \(Z\) supported on the closed
subset \(E\) is canonically such a sheaf.  Thus \(g_*\) and its derived
functors are defined for \(\mathcal T_{j,k}\), although there is no
holomorphic map from its thickening to \(U\).  For a coherent layer
\(\mathcal A_j\), this is the usual coherent analytic direct image.
The finite-neighborhood assertion we will prove is the following.

\begin{proposition}[All finite lifting orders]\label{prop:all-orders}
For every local root chart \eqref{eq:local-root}, every integer
\(j\), and every \(k\geq1\), the map
\begin{equation}\label{eq:all-order-surjection}
 g_*\mathcal T_{j,k+1}\longrightarrow g_*\mathcal T_{j,k}
\end{equation}
is an epimorphism of sheaves of complex vector spaces on \(U\).
The neighborhoods used to lift a particular germ may depend on \(j\)
and \(k\).
\end{proposition}

The proof is in Section~\ref{sec:lifting}.  To see what geometry it
requires, consider the exact sequence
\begin{equation}\label{eq:current-truncation-sequence}
 0\longrightarrow\mathcal A_{j+k}\longrightarrow\mathcal T_{j,k+1}
 \longrightarrow\mathcal T_{j,k}\longrightarrow0.
\end{equation}
Its connecting map takes a section of \(g_*\mathcal T_{j,k}\) to an
obstruction in \(R^1g_*\mathcal A_{j+k}\).  In order \(k\), choosing
\(j=-a-k\) puts that obstruction in the fixed sheaf
\(R^1g_*\mathcal A_{-a}\).  We next construct a split insertion of this
sheaf into the first cohomology of an SNC dualizing sheaf.  The Hodge
argument will apply there, and the induction will return from these
negative degrees to all integer degrees.

\subsection{A canonical root and its resolved support}

Take a second full normalized cover, this time adjoining a \(q\)th root
of \(\pi^*s\) as a meromorphic \(q\)-pluricanonical form.  The local
construction in Lemma~\ref{lem:floor-torsion-free} also proves its
existence on the normal analytic \(Z\): if its coefficient in a
meromorphic canonical generator is \(a_0/b_0\), normalize the finite
reduced monic algebra
\[
 A[w]/(w^q-a_0b_0^{q-1})
 \subset \operatorname{Frac}(A)[t]/(t^q-a_0/b_0).
\]
The total meromorphic algebra is separable and its components all dominate.
A change of generator identifies the total meromorphic algebras by
multiplying the root by a meromorphic unit.  The displayed finite monic
orders need not coincide under this identification, but their integral
closures do: each is the integral closure of the base ring in that same
total algebra.  These closures therefore glue.  The forms
\(t\eta\) glue to the tautological meromorphic top form \(\tau\).
For the total meromorphic field \(K\) of a normal base component, the
full \(\mu_q\)-action has invariant algebra \(K\) and root-character
space \(Kt\); the corresponding meromorphic form space is \(K\tau\).
The invariant subalgebra of the normalized holomorphic algebra is the
base ring \(A\), by normality.  These statements are read componentwise
on a disconnected base.
The earlier \(\mu_r\)-action lifts functorially, preserving the pulled-back
pluriform, and commutes with \(\mu_q\).  Choosing one component instead
would in general destroy these assertions.

Resolve this second cover and principalize the pulled ideal of \(E\),
equivariantly for these finite actions.  We use the analytic
smooth-functorial resolution and principalization of
\cite[Theorems~1.1.11 and~1.1.13]{Temkin2017}, together with its nonembedded
analytic construction \cite[Sections~5.2.2 and~5.3.2]{TemkinNonembedded}.
We retain the final indexed complete boundary in the construction of
\(F_{\mathrm{princ}}\): after the blowup of the ideal, boundary
desingularization makes the final strictly monomial support a union of
components of this SNC boundary.  Its closed indexed intersections are
regular by \cite[Lemma~2.1.10]{Temkin2017}.  Splitting the disjoint
connected components of each boundary component near the compact fiber
gives globally smooth component labels; their intersections remain
smooth, and the finite group may permute the labels.  On a neighborhood of a compact
parameter fiber only finitely many stages of the locally finite blowup
hypersequence occur; a finite-group invariant shrink preserves equivariance.
Write the composite as \(\rho:\widehat Z\to Z\), with \(\widehat Z\)
smooth, and put
\begin{equation}\label{eq:resolved-boundary}
 D_E=\rho^*E,\qquad H=(D_E)_{\red},\qquad
 \rho_H:H\to E,\qquad h=g\rho_H.
\end{equation}
Here \(D_E\) is the scheme inverse Cartier divisor and \(H\) is its
globally simple SNC reduction.  The tautological form satisfies
\(\tau^{\otimes q}=(\pi\rho)^*s\) with differential pullback understood.

These constructions commute with restrictions to opens and ordinary
products by a complex manifold.  For normalization, the product of the
normal normalization with a manifold is finite, normal, and bimeromorphic
to the reduced product, hence is its normalization by uniqueness;
normality of these analytic products is part of
\cite[Part~C, Theorem~4(b)]{Houzel1961}.
The resolution is functorial for smooth morphisms, including products.
We make no assertion about normalization or resolution under a ramified
base change.

We will also need the K\"ahler property near the resolved support.  A finite
analytic map is projective locally: generators of its finite algebra embed
it in a relative projective space, and the affine coordinate \(1\) makes
the trivial line relatively ample.  The finite maps and the finitely many
blowups above are therefore projective near the preimage of a compact set.
Patch local positive metrics of a relatively ample line by partitions
pulled from the base.  Their curvature remains positive on vertical tangent
directions.  A sufficiently large multiple of the pulled-back K\"ahler
form of \(W\) makes it positive in all directions near the compact
preimage; compactness bounds the mixed terms uniformly.  This gives a
closed K\"ahler form on an ambient neighborhood \(\mathcal N\) of the
compact resolved fiber under consideration.  Properness of \(h:H\to U\)
lets us shrink \(U\) so that all of \(H\) above the smaller \(U\) lies in
\(\mathcal N\): remove the closed image of \(H\setminus\mathcal N\).
We make this shrink in each root chart.  Every closed intersection of
components of \(H\) is then smooth and proper over \(U\), and inherits
one global relative K\"ahler form there.  A connected component can split
after a further restriction of the base, so we make the indexing choice
only after these neighborhood restrictions.

\subsection{Finite indexing near a parameter fiber}

The filtered direct-image construction will use the closed intersections
of globally indexed smooth components.  The following lemma makes that
indexing finite after a base restriction, even when some strata map only
to special parameter loci.

\begin{lemma}[Components near a compact fiber]\label{lem:finite-components}
Let \(h:H\to T\) be a proper holomorphic map to a complex manifold, and
fix \(t_0\in T\).  Suppose that near \(h^{-1}(t_0)\) the reduced support
\(H\) has a locally finite family of closed smooth labels \(D_\lambda\)
which locally are its distinct SNC branches, and that every closed
intersection of these labels is smooth.  After restricting to a suitable
open neighborhood \(U\) of \(t_0\), the support has finitely many globally
smooth indexed irreducible components.  Every closed intersection of the
new components has finitely many connected components, all smooth and
proper over \(U\).  A global relative K\"ahler form already given on the
old strata restricts to such a form on the new ones.  The open \(U\) can
be chosen inside any previously prescribed neighborhood of \(t_0\).
\end{lemma}
\begin{proof}
Local finiteness and compactness give an open neighborhood
\(\mathcal N_1\) of the central fiber which meets only finitely many
labels.  Remove from the base the closed image of
\(H\setminus\mathcal N_1\); properness then puts the entire restricted
support in \(\mathcal N_1\).  There are now finitely many closed indexed
intersections to consider.  The connected components of each form a
locally finite family in \(H\): the intersection is closed and smooth,
so near one of its points a small smooth neighborhood meets only its own
component, and near a point outside it a neighborhood avoids it.
Compactness, followed by the same properness argument, gives a second
neighborhood \(\mathcal N_2\) meeting only finitely many such components
and a base restriction on which the whole support lies in
\(\mathcal N_2\).

Let \(C_1,\ldots,C_N\) be the whole connected components of the
pre-restriction closed intersections which meet \(\mathcal N_2\).
We retain the whole components, not just their intersections with
\(\mathcal N_2\).  Each \(C_a\) is closed in a closed stratum and is a
connected complex manifold.  Thus \(f_a=h|_{C_a}\) is proper over the
preceding base and \(C_a\) is irreducible.  Their restrictions cover
every closed intersection after the current restriction, including the
singleton intersections.

Put \(S_a=f_a(C_a)\) with its reduced structure.  Remmert's proper
mapping theorem makes \(S_a\) closed analytic
\cite[Section~7, no.~1, Satz~1, p.~60]{Grauert1960ProperMapping}, and it
is irreducible because \(C_a\) is.  Inside any prescribed open coordinate
neighborhood \(V\ni t_0\) contained in the preceding restrictions,
take the locally finite irreducible decompositions of all \(S_a\cap V\).
Let \(B\) be the union of those components which do not contain \(t_0\).
A union of a subfamily of the locally finite irreducible components of
an analytic set is closed analytic.  Since the list of \(a\)'s is finite,
\(B\) is closed analytic in \(V\) and misses \(t_0\).  Set \(U=V\setminus B\).

This removal does not split any retained irreducible component \(A\) of
an \(S_a\cap V\).  Its regular locus is connected, and its intersection
with \(B\) is a proper analytic subset: otherwise closedness and density
would imply \(A\subset B\), contrary to \(t_0\in A\setminus B\).
The complement of a proper analytic subset in a connected complex
manifold is connected.  One may see this by perturbing a path, in a
finite chain of coordinate balls, off a subset of positive complex
codimension.  Hence \(\operatorname{Reg}(A)\setminus B\) is connected
and \(A\cap U\) is irreducible.  Only finitely many components of
\(S_a\cap V\) contain \(t_0\), by local finiteness.  Their restrictions
give a finite cover of \(S_a\cap U\) by irreducible closed analytic
subsets, all containing \(t_0\).  Therefore every irreducible component
of \(S_a\cap U\) contains \(t_0\); if \(S_a\) misses \(t_0\), its
restriction is empty.

Let \(W\) be a connected component of \(f_a^{-1}(U)\).  It is open in
the connected smooth \(C_a\), is closed in \(f_a^{-1}(U)\), and is
smooth and irreducible.  Its map to \(U\) is proper.  Put
\(r_a=\dim S_a\).  The maximal-rank locus of \(f_a\) is dense in \(C_a\):
the generic rank theorem gives maximal rank \(r_a\), and the lower-rank
locus is a proper analytic subset defined by the maximal minors.
The nonempty open \(W\) meets the maximal-rank locus.  Its proper image is consequently
an irreducible closed analytic subset of \(S_a\cap U\) of dimension
\(r_a\).  Every component of \(S_a\cap U\) has this dimension because
\(S_a\) is irreducible and hence pure-dimensional.  The image is therefore
an irreducible component, since a proper analytic subset of such a
component has smaller dimension.
It therefore contains \(t_0\), and \(W\) meets \(f_a^{-1}(t_0)\).

The compact analytic fiber has finitely many connected components:
on its reduction, irreducible components are locally finite and
compactness makes their number finite.  Each
\(W\cap f_a^{-1}(t_0)\) is a nonempty open-and-closed subset of this
fiber.  Different \(W\)'s give disjoint such subsets, so
\begin{equation}\label{eq:finite-fiber-components}
 \#\pi_0(f_a^{-1}(U))
 \leq \#\pi_0\bigl((f_a^{-1}(t_0))_{\red}\bigr)<\infty.
\end{equation}
This includes zero-dimensional and other vertical images and permits
nonreduced scheme fibers.

Reindex the singleton strata by these finitely many components \(W\).
An intersection of the new labels is an open-and-closed part of the
restriction of the corresponding old closed intersection, hence is a
union of some of its finitely many connected components.  It remains
smooth and proper, and the relative form restricts.  Distinct new pieces
from one old label are disjoint, so the local SNC branches and their
ordering are unchanged.  This proves the assertion.
\end{proof}

Apply Lemma~\ref{lem:finite-components} after the preceding K\"ahler
neighborhood and properness restrictions, to all regular labels and their
closed intersections.  The resulting local support has the finite global
component indexing used in Section~\ref{sec:hodge}.  The final open need
not be Stein or a polydisc: the roots and covers have already been
constructed, and their remaining local uses only restrict the existing
coordinate charts.

\subsection{A residue insertion that retains all cohomology}

\begin{proposition}[Split residue insertion]\label{prop:residue-insertion}
For the data \eqref{eq:cover-layers}--\eqref{eq:resolved-boundary},
multiplication by \(\tau\) gives a morphism
\[
 \rho^*\OO_Z(aE)\longrightarrow\omega_{\widehat Z}(H),
\]
and residue gives \(\rho_H^*\mathcal A_{-a}\to\omega_H\).  The resulting
map in the derived category of \(E\) has a retraction in the
\(\tau\)-character.  Consequently, for every \(j\), and in particular
for \(j=1\), it gives a split injection
\begin{equation}\label{eq:residue-injection}
 \tau_*:R^jg_*\mathcal A_{-a}\lhook\joinrel\longrightarrow R^jh_*\omega_H.
\end{equation}
The construction is natural under the power-compatible root comparisons
of Lemma~\ref{lem:root-neighborhood}.  Under an ordinary product with a
complex manifold \(B_0\), it is the pullback construction for canonical
sheaves relative to \(B_0\).  For absolute canonical sheaves the insertion
and its retraction are tensored with \(\omega_{B_0}\); in particular the
product insertion has source \(\mathcal A_{-a}\boxtimes\omega_{B_0}\)
and target \(R(\rho_H\times\id)_*\omega_{H\times B_0}\).
\end{proposition}
\begin{proof}
Let \(v\) be a local meromorphic frame of \(\OO_Z(aE)\).  By
\eqref{eq:cover-line} and \(aq=r\), the \(q\)-pluriform
\(v^q\pi^*s\) is, up to a holomorphic unit, the pullback of a local frame
\(\ell_V\) of \(L\).  On a log resolution of \((V,B,G)\), write \(b_F\)
for a crepant boundary coefficient and
\(m_F=\ord_F(s/\ell_V)\).  Then \(m_F\geq0\), \(b_F\leq1\), and
\(b_F=1\) implies \(m_F>0\): the dlt pair is klt off \(S\), and the
section vanishes exactly on \(S\).  Thus the local density
\[
 |\ell_V|^{2/q}|s/\ell_V|^{2\epsilon}
\]
has exponent \(-b_F+\epsilon m_F>-1\) at each prime for every
\(\epsilon>0\), and is locally integrable.  Change of variables under
the proper generically finite map \(\pi\rho\) preserves this integrability.
For \(\alpha=(\rho^*v)\tau\), let \(k_F\) be its integral order at a
prime upstairs and \(n_F\) the order of the pulled section ratio.
Integrability says \(k_F+\epsilon n_F>-1\) for all \(\epsilon>0\).
The order \(n_F\) is positive exactly on \(H\).  Hence \(k_F\geq-1\)
on \(H\) and \(k_F\geq0\) elsewhere.  This proves the logarithmic map.
Cartier adjunction on the reduced SNC divisor gives its residue map on
\(H\), with \(\rho_H^*\mathcal A_{-a}=\rho^*\OO_Z(aE)|_H\).

Let \(\chi\) be the \(\tau\)-character of the second cover.  At a general
point of \(E\) over \(S_i\), differential pullback gives
\[
 \ord_E(\pi^*s)
 =\frac r{d_i}(d_i-q)+q\left(\frac r{d_i}-1\right)
 =r-q=q(a-1).
\]
The second root is unramified there and \(\ord_E\tau=a-1\).  At any
other codimension-one prime the first cover is \'etale and the original
boundary coefficient is \(\beta\in[0,1)\).  For
\(e=q/\gcd(q,q\beta)\), the second-cover order is
\(e(1-\beta)-1\in\{0,\ldots,e-1\}\), by
\eqref{eq:canonical-root-order}.  A meromorphic form in character
\(\chi\) is \(f\tau\) for an invariant meromorphic \(f\) descended
to \(Z\).  The orders force \(\ord_E f\geq-(a-1)\) and nonnegative
orders at all other primes.  Normality extends it as a section of the
indicated invertible line.  Conversely, the logarithmic map multiplied
by the equation of \(E\) has no pole because \(D_E\geq H\).  We obtain
the exact sheaf identities
\begin{align}
 (\rho_*\omega_{\widehat Z})_\chi
   &=\OO_Z((a-1)E)\tau,\label{eq:residue-character-one}\\
 (\rho_*\omega_{\widehat Z}(D_E))_\chi
   &=\OO_Z(aE)\tau.\label{eq:residue-character-two}
\end{align}
The second is projection formula for the invertible line \(\OO_Z(E)\).

Canonical torsion-freeness
\cite[Theorem~2.9 and Proposition~2.11]{Fujino2023AnalyticVanishing}
applies locally near \(E\) to each dominating resolved component.  More
precisely, around a compact fiber choose the K\"ahler neighborhood just
constructed.  Properness of \(\rho\) supplies a smaller normal target
neighborhood whose entire inverse image lies in it, by removing the
closed image of its complement.  After a base shrink this neighborhood
contains all of \(E\) in the root chart.  The restricted morphism is
proper with K\"ahler smooth source, so the cited theorem applies there.
These neighborhoods cover \(E\).  On each, the higher direct images for
the generically finite resolution vanish because they vanish generically.
Finite pushforward is exact.  Thus on their union \(Z^\circ\subset Z\),
an open neighborhood of \(E\),
\[
 (R^i\rho_*\omega_{\widehat Z})|_{Z^\circ}=0\quad(i>0).
\]
Projection formula gives the same restricted vanishing for
\(\omega_{\widehat Z}(D_E)\).

Put \(A_E=\mathcal A_{-a}=\OO_E(aE)\).  Cartier adjunction for the
possibly nonreduced divisor \(D_E\) is the exact sequence
\[
 0\longrightarrow\omega_{\widehat Z}\longrightarrow
 \omega_{\widehat Z}(D_E)\longrightarrow\omega_{D_E}\longrightarrow0.
\]
Since \(D_E\) is the scheme inverse image of \(E\), it has a morphism
\(\rho_D:D_E\to E\).  The preceding vanishings and character identities
identify
\begin{equation}\label{eq:residue-quotient}
 \bigl(R(\rho_D)_*\omega_{D_E}\bigr)_\chi\simeq A_E\tau[0]
 \quad\text{in }D(\OO_E).
\end{equation}
One can see this first after exact closed pushforward from \(E\) to
\(Z^\circ\), where it
is the quotient of \eqref{eq:residue-character-two} by
\eqref{eq:residue-character-one}.  Closed pushforward detects its
cohomology sheaves, and the canonical truncation of a complex concentrated
in degree zero gives \eqref{eq:residue-quotient} on \(E\).  No derived
full-faithfulness assertion is needed.

Because \(A_E\) is invertible, its pullback followed by the residue map
defines \(A_E[0]\to R(\rho_H)_*\omega_H\).  The inclusion
\(\omega_H\hookrightarrow\omega_{D_E}\), followed by the character
projection and \eqref{eq:residue-quotient}, gives a map back to
\(A_E\tau[0]\).  On degree zero their composition sends \(f\) to
\(f\tau\) in the quotient of the two character lines.  Under the chosen
identification it is the identity.  An endomorphism of the sheaf
\(A_E[0]\) in degree zero is determined by its sheaf map, so this is a
derived retraction.  Applying \(Rg_*\) proves the split injections
\eqref{eq:residue-injection}.  All maps were defined by the tautological
form, the actual divisor ideal, and residue, so they respect root
comparisons.  Under an ordinary product the tautological form is relative
to the new factor; wedging with a local frame of its canonical line gives
the absolute map and the stated \(\omega_{B_0}\) factor.
\end{proof}

The derived retraction is stronger than an injection only on functions or
at general fibers.  It is this strength that allows the later argument
to retain obstruction classes supported on special parameters.

\subsection{One compact graph on a parameter cover}

We next prepare the geometric setting for a global symbol test.  For any
chosen parameter \(t_*\), the construction gives a projective parameter
cover unbranched over \(t_*\) and one compact SNC graph over the entire
cover.  Global vanishing can then be tested back at \(t_*\) without
discarding a class supported there.  The cover may change with \(t_*\).
We form and resolve the covers before imposing the graph equations; this
avoids assuming that normalization commutes with ramified base change.

\begin{proposition}[Compact transverse graph]\label{prop:global-graph}
Assume \(b>0\), and fix any \(t_*\in T=\PP^b\).  There is a coordinate-power
map \(\psi:T'=\PP^b\to T\) of degree \(r\) in each coordinate, with
\(R=\OO_{T'}(1)\) and an actual identity \(R^r\simeq\psi^*\OO_T(1)\),
which is unbranched over \(t_*\), and the following data.  On a neighborhood
of the compact graph \(\Xi=S\times_T T'\subset V\times T'\) there are
the two full root covers and a projective resolution with smooth ambient
source \(\mathscr W\).  The graph cut \(H'\) inside the reduced inverse
section divisor is compact and reduced SNC of pure dimension \(\dim V-1\),
with finitely many globally smooth components.  The projection
\(p':\mathscr W\to T'\) is a submersion near \(H'\), and every nonempty
closed component intersection is smooth and proper over \(T'\), carrying
the restriction of one K\"ahler form on a neighborhood of \(H'\).

Locally near every compact graph slice over \(w\in T'\), the resolved
ambient data, including their projection and full tautological divisor
ideal, are isomorphic to an open restriction of the ordinary product
\(\widehat Z\times U'\) for a local root chart.  There the cut is
\(H^\flat=H\times_U U'\), and its actual dualizing line is
\begin{equation}\label{eq:graph-adjunction}
 \omega_{H^\flat}\simeq\operatorname{pr}_H^*\omega_H\otimes
 \operatorname{pr}_{U'}^*(\omega_{U'}\otimes\psi^*\omega_U^{-1}).
\end{equation}
\end{proposition}
\begin{proof}
Choose finitely many pairs of parameter opens \(U_i'\Subset U_i\), with
the \(U_i'\) covering \(f(S)\), such that a local root chart and its
proper map \(h_i:H_i\to U_i\) are defined on \(U_i\).  Each
\(h_i^{-1}(\overline{U_i'})\) is compact.  Local finiteness therefore
leaves only finitely many smooth component intersections meeting these
compact sets.  For each such stratum \(C\) and each subset of
variable hyperplanes in \(T\), consider the incidence
\[
 \{(z,H_1,\ldots,H_v):h(z)\in H_1\cap\cdots\cap H_v\}.
\]
It is smooth, since varying each hyperplane independently supplies its
normal direction.  Sard's theorem for the projection to the hyperplane
parameter space gives transversality for a general tuple; a countable
atlas handles any noncompact stratum.  Choose a tuple satisfying all
the finitely many conditions and also the open conditions that it forms
a coordinate basis and that none of its hyperplanes contains \(t_*\).
Use those coordinates for
\[
 \psi([w_0:\cdots:w_b])=[w_0^r:\cdots:w_b^r].
\]
Its derivative at a point has image the tangent space to the intersection
of exactly the coordinate hyperplanes corresponding to vanishing
coordinates, or is invertible when none vanish.  The chosen transversality
therefore gives, at \(h(z)=\psi(w)\),
\begin{equation}\label{eq:graph-transversality}
 dh(T_zC)+d\psi(T_wT')=T_{h(z)}T
\end{equation}
for every smooth SNC stratum.  The map is unbranched over \(t_*\).

In local coordinates \(t_i\) on \(U\), extend the coordinate functions
of \(h\) holomorphically off \(H\) to functions \(G_i\) on local opens
of \(\widehat Z\).  In \(\widehat Z\times U'\), the equations
\(\psi_i(w)-G_i=0\) have surjective differential on each \(H\)-stratum
by \eqref{eq:graph-transversality}.  They cut a smooth submanifold of
codimension \(b\) transverse to the SNC divisor.  The cut
\(H^\flat=H\times_U U'\) is therefore reduced SNC of pure dimension
\(\dim V-1\), and has codimension \(b+1\) in the product ambient.
Complete-intersection adjunction gives \eqref{eq:graph-adjunction}.
In coordinates its last factor has the wedge frame \(dw/dt\); this means
a frame of \(\omega_{U'}\otimes\psi^*\omega_U^{-1}\), not an inverse
Jacobian.

The graph \(\Xi\) is a compact closed analytic subspace.  Its actual
line identity is
\[
 \operatorname{pr}_V^*L|_\Xi\simeq
 \operatorname{pr}_{T'}^*R^r|_\Xi.
\]
Lemma~\ref{lem:root-neighborhood} gives a root \(P_*\) of
\(\operatorname{pr}_V^*L\) on a neighborhood of \(\Xi\), identified
with \(\operatorname{pr}_{T'}^*R\) along \(\Xi\).  Form the same two
full normalized covers there.  For the second cover use coefficients
of \(s\) in meromorphic canonical generators from the first factor
\(V\).  This defines its relative tautological form without requiring a
Cartier relative canonical line on the singular cover; changing to a
generator after differential pullback changes the coefficient by a
\(q\)th power and identifies the total root algebra and its normalization.
Resolve and principalize the full tautological ideal functorially on this neighborhood, before
cutting by the graph.

Near \(w_0\), choose a frame \(\bar p\) of \(R\) whose \(r\)th power
is \(\psi^*\ell\) for a downstairs frame \(\ell\); a local root of a
unit gives such a choice.  Along the compact graph slice, compare the
pullback of \(P_1\) with \(P_*\) by sending the boundary frame \(p_1\)
to \(\bar p\).  Lemma~\ref{lem:root-neighborhood} extends this to an
ambient root isomorphism as a germ.  On the graph, its ratio with the
prescribed comparison is a \(\mu_r\)-section equal to one on the compact
slice, hence equal to one on an open graph neighborhood of that slice.
Properness of \(\Xi\to T'\) permits a base shrink on which the ambient
comparison is defined and agrees with the prescribed one along the entire
graph: remove the closed image of the complement of that neighborhood.
This argument also applies to a nonreduced graph.  The comparison identifies
the first cover with the local ordinary product cover, and product
normality identifies the second cover as well.  Apply the same fixed
smooth-functorial principalization to the same full ideal and ordered
boundary data on both sides; functoriality identifies the results.  These are isomorphisms of resolved
ambient germs over the product base preserving the whole divisor ideal,
not only isomorphisms of their supports.

Let \(H'\) be the graph cut inside the reduced inverse section divisor
on this global resolution.  Locally it is exactly \(H^\flat\) above,
so it is reduced SNC with the asserted pure dimension.  It is compact,
being over the compact \(\Xi\) by proper maps.  The ambient projection
is a submersion near it by the local product comparison.  The transverse
cuts of the globally smooth inverse-divisor components are smooth; split
them into their disjoint connected components.  Compactness and local
finiteness leave only finitely many.  Their closed intersections are
smooth and compact, hence proper over \(T'\).  Finally, the relative
metric construction above, applied over the K\"ahler product neighborhood
of \(\Xi\), gives one K\"ahler form near \(H'\).  Its restrictions supply
all the required relative K\"ahler forms.
\end{proof}

The local and global graph supports in this proposition are locally a
reduced SNC divisor in a smooth complete intersection, with submersive
ambient projection.  The next section develops the filtered direct-image
statement for precisely that geometry.
